\appendix

\section{Appendix A}

Mathematical Foundations

Appendix A is a formula index for the main text. It repeats notation already defined in the chapters and introduces no additional operators.

A.1 Unified Constraint Operator

\[ \mathbb{L}:\mathcal{S}\rightarrow\mathcal{C}. \]

A.2 Observable Projection

\[ L_i=\pi_i(\mathbb{L}), \qquad i=1,\ldots,n. \]

A.3 Projection Vector

\[ \mathbf{L}=(L_T,L_F,L_C,L_S,L_I,L_E,L_G,L_V). \]

A.4 Optional Scalar Index

\[ L=F(\mathbf{L}), \qquad F:[0,1]^n\rightarrow[0,1]. \]

\[ L\approx\sum_i w_i L_i, \qquad \sum_i w_i=1. \]

A.5 Constraint Metric

\[ ds_{\mathbb{L}}^2=g_{ij}(\mathbf{L})\,dL_i\,dL_j. \]

A.6 Constraint Dynamics

\[ \frac{D\mathbb{L}}{Dt}=\mathcal{F}(\mathbb{L},\Psi,t). \]

A.7 Coupling Matrix

\[ C_{ij}=\frac{\partial \pi_i(\mathbb{L})}{\partial \pi_j(\mathbb{L})}. \]

A.8 Global Coupling

\[ \Gamma=\sum_{i\neq j}|C_{ij}|. \]

A.9 Viability Functional

\[ \mathcal{V}=\Phi(\mathbb{L}) \quad \text{or} \quad \mathcal{V}=\Phi(\mathbf{L}). \]

\section{Appendix B}

Dimensionless Normalization

Different observables possess different physical units and uncertainty structures. UCT therefore uses dimensionless projected components, typically scaled to \([0,1]\), while recognizing that normalization is domain-specific.

\[ L_i=\frac{x_i-x_i^{\mathrm{ref}}}{x_i^{\mathrm{crit}}-x_i^{\mathrm{ref}}}. \]

When a bounded index is required, clipping may be applied:

\[ L_i\leftarrow \min(1,\max(0,L_i)). \]

Alternatively, a monotonic transform may be used:

\[ L_i=\sigma_i(x_i). \]

The reference value \(x_i^{\mathrm{ref}}\), critical value \(x_i^{\mathrm{crit}}\), and transform \(\sigma_i\) must be chosen and reported for each projection domain. Spatial scaling, temporal scaling, energy scaling, and information scaling can all be incorporated through projection operators.

\section{Appendix C}

Representative Observational Data Sources

These examples are illustrative. A validated application must specify data provenance, preprocessing, uncertainty treatment, and the projection operator used for each domain.

Table~\ref{tab:data_sources} summarizes representative data sources for each projection.

\begin{table}[tb]
\caption{\label{tab:data_sources}Representative observational data sources.}
\footnotesize
\begin{ruledtabular}
\begin{tabular}{lll}
Projection & Observable & Example datasets \\
Thermal & TOA imbalance; temperature anomaly & CERES; ERA5; NOAA \\
Flow & Wind; vorticity; AMOC proxies & ERA5; ECMWF; Argo; RAPID \\
Chemical & CO$_2$; CH$_4$; ocean pH & NOAA; Mauna Loa; SOCAT \\
Structural & Land; ice; infrastructure & MODIS; Landsat; Sentinel; GRACE \\
Informational & Biodiversity; complexity & GBIF; IUCN; LPD \\
Energetic & Availability; efficiency & IEA; Energy statistics \\
Geophysical & Boundary conditions & GRACE; Geodesy \\
Viability & Population; infrastructure; capacity & UN; World Bank; Statistics \\
\end{tabular}
\end{ruledtabular}
\end{table}
